\begin{tikzpicture}[
  font=\scriptsize,
  elemen/.style={draw, minimum height=0.95cm, text width=2.0cm, align=center,
                 font=\scriptsize\bfseries},
  panah/.style={-{Stealth[length=2mm]}, thick},
  wujud/.style={-{Triangle[open, length=2mm, width=2mm]}, thick, dashed},
  kaitan/.style={thin, gray!80!black}
]
  % Plateau
  \node[elemen, fill=gray!35] (awal) at (0,0)
    {{\tiny\mdseries\itshape plateau}\\Kondisi awal 2026};
  \node[elemen, fill=blue!8] (t1) at (4.7,0)
    {{\tiny\mdseries\itshape plateau}\\T1 Maret 2028};
  \node[elemen, fill=green!10] (t2) at (9.4,0)
    {{\tiny\mdseries\itshape plateau}\\T2 Maret 2029};
  \draw[panah] (awal) -- (t1);
  \draw[panah] (t1) -- (t2);

  % Gap
  \node[elemen, fill=red!22, text width=2.15cm] (g1) at (2.35,1.5)
    {{\tiny\mdseries\itshape gap}\\Data pasien ganda, akses tak terkendali};
  \node[elemen, fill=red!22, text width=2.15cm] (g2) at (7.05,1.5)
    {{\tiny\mdseries\itshape gap}\\Rekam medis terpisah per rumah sakit};
  \draw[kaitan] (g1.south west) -- (awal.north);
  \draw[kaitan] (g1.south east) -- (t1.north);
  \draw[kaitan] (g2.south west) -- (t1.north);
  \draw[kaitan] (g2.south east) -- (t2.north);

  % Deliverable
  \node[elemen, fill=blue!8] (d1) at (3.1,-1.6)
    {{\tiny\mdseries\itshape deliverable}\\Direktori akses terpusat};
  \node[elemen, fill=blue!8] (d2) at (6.1,-1.6)
    {{\tiny\mdseries\itshape deliverable}\\Indeks pasien tunggal};
  \draw[wujud] (d1.north) -- (t1.south west);
  \draw[wujud] (d2.north) -- (t1.south east);

  % Work package dan implementation event
  \node[elemen, fill=orange!12, rounded corners=6pt] (w1) at (3.1,-3.1)
    {{\tiny\mdseries\itshape work package}\\P1 Kontrol akses};
  \node[elemen, fill=orange!12, rounded corners=6pt] (w2) at (6.1,-3.1)
    {{\tiny\mdseries\itshape work package}\\P3 Identitas pasien};
  \node[elemen, fill=praditagreen!15, rounded corners=2pt] (k1) at (8.9,-3.1)
    {{\tiny\mdseries\itshape implementation event}\\K1 Anggaran 2027 disahkan};
  \draw[wujud] (w1) -- (d1);
  \draw[wujud] (w2) -- (d2);
  \draw[panah] (w1) -- (w2);
  \draw[panah] (k1) -- (w2);

  \node[anchor=west, text width=11.4cm, font=\tiny] at (-1.15,-4.25)
    {\textbf{Panah tebal}: \textit{triggering}, urutan waktu.
     \textbf{Garis putus-putus berujung segitiga}: \textit{realization}.
     \textbf{Garis tipis}: \textit{association} antara \textit{gap} dan dua
     \textit{plateau} yang dibandingkannya.};
\end{tikzpicture}
